% Generated draft: requires mathematical and source review.
\documentclass[11pt]{article}
\usepackage[margin=25mm]{geometry}
\usepackage{fontspec}
\setmainfont{Latin Modern Roman}
\usepackage{amsmath,amssymb,mathtools,unicode-math}
\setmathfont{Latin Modern Math}
\usepackage{enumitem,needspace,xurl,hyperref}
\setlength{\parindent}{0pt}
\setlength{\parskip}{5pt}
\setlength{\emergencystretch}{3em}
\newcommand{\N}{\mathbb{N}}
\newcommand{\Z}{\mathbb{Z}}
\newcommand{\Q}{\mathbb{Q}}
\newcommand{\R}{\mathbb{R}}
\newcommand{\C}{\mathbb{C}}
\newcommand{\F}{\mathbb{F}}
\newcommand{\A}{\mathbb{A}}
\newcommand{\PP}{\mathbb{P}}
\newcommand{\E}{\mathbb{E}}
\newcommand{\eps}{\varepsilon}
\newcommand{\dd}{\,\mathrm{d}}
\begin{document}
\section*{UnsolvedMath ID: 30001309}
\textbf{Source problem code:} catalogue entry
\subsection{Definitions and mathematical statement}
A potential flow is a flow whose velocity field admits a velocity potential: writing \(\Omega(t)\) for the fluid region, one has \(\boldsymbol{u}(x,t)=\nabla\phi(x,t)\) in that region. A free surface is an unknown part \(\Gamma(t)\) of the boundary of \(\Omega(t)\), whose evolution is part of the problem. These descriptions do not specify a particular governing system or its boundary conditions.

The supplied question is: \emph{Is the free-surface potential-flow problem in hydrodynamics integrable?}

The record does not define \emph{integrable} or specify the spatial dimension, fluid geometry, depth, governing equations, physical assumptions, forcing, surface tension, boundary conditions, or class of admissible solutions. In particular, it does not specify whether integrability means a suitable infinite-dimensional analogue of Liouville integrability, an inverse-scattering formulation, or another property. Consequently, the supplied question cannot yet be expressed as a uniquely specified mathematical proposition with definite hypotheses and quantifiers. A precise formulation requires the intended free-surface evolution system and the intended criterion of integrability.

The supplied dated curation labels the problem partially solved and states that the complete-integrability question remains unresolved, that nonintegrability has been proved for a compact one-dimensional Zakharov reduction, and that special regimes admit extra integrals and exact solutions. The reduction and special regimes are not defined in the record, so these qualifications do not identify the precise target system or settle the stated question here.
\subsection{Short English statement}
The question asks whether free-surface potential flow has an integrable mathematical structure. The supplied record leaves both the fluid model and the meaning of integrability unspecified. Its partially solved label and claims about reductions and special regimes are supplied dated assessments requiring independent review.
\subsection{Sources}
\begin{itemize}
\item Supplied local UnsolvedMath registry record, permanent ID \(30001309\), title \emph{Integrability of Free-Surface Potential Flow}, Oberwolfach Reports Open Problems collection, source problem code \texttt{catalogue entry}. The record gives proposed year 2009 and identifies no individual proposer.
\item Supplied source citation and literature link labeled \emph{Hydrodynamics report}: \emph{Mathematical Aspects of Hydrodynamics} (2010), Oberwolfach Reports, pp.\ 1921--1942, DOI 10.4171/owr/2009/34. \url{https://doi.org/10.4171/owr/2009/34}. The record attributes the question to workshop open-problem discussions documented in Oberwolfach Reports.
\item Supplied literature link labeled \emph{Nonintegrability result}: \url{https://doi.org/10.7868/S0370274X13130109}. No authors, title, hypotheses, or theorem statement are supplied.
\item Supplied literature link labeled \emph{Free-surface integrability study}: \url{https://arxiv.org/abs/1604.04778}. No authors, title, or precise model are supplied.
\item Literature assessment date: \texttt{2026-08-21}. This is supplied dated curation, not an independently verified assessment. It labels the problem partially solved, reports an unresolved complete-integrability question and nonintegrability for a compact one-dimensional Zakharov reduction, and reports extra integrals and exact solutions in special regimes. The source expressly requests independent verification of time-sensitive status and cited claims.
\end{itemize}
\end{document}
